\documentclass[11pt,a4paper]{article}

%% PACHETE MATEMATICE
\usepackage{amsmath,amssymb,amsfonts}
\usepackage{amsthm}
\usepackage{mathpazo}
\usepackage{eulervm}
\usepackage{enumerate}
\usepackage{color}
\usepackage{graphicx}
\usepackage[all]{xy}
\usepackage{xcolor}
\usepackage{framed}
\usepackage[unicode]{hyperref}
\setcounter{MaxMatrixCols}{20}
\usepackage{multicol}
%\usepackage[romanian]{babel}


%% ASPECT
\linespread{1.05}
\usepackage[T1]{fontenc}
\usepackage[utf8x]{inputenc}
\usepackage[margin=2cm]{geometry}
% \usepackage[color]{showkeys}
\usepackage{anyfontsize}
\usepackage{titlesec}
\titleformat*{\section}{\large\bfseries}

\usepackage{fancyhdr}
\pagestyle{fancy}
\rhead{\small \color{gray} Notițe de Adrian Manea}
\lhead{\small \color{gray} BCD-Aero, 2017---2018, M. Olteanu}


\usepackage[autostyle = false, style = german]{csquotes}
\usepackage[romanian]{babel}
\newcommand\qq{\enquote}

%% COMENZILE MELE
\renewcommand*{\proofname}{Dem.:}
\newcommand\bb{\mathbb}
\newcommand\dr{\mathrm}
\newcommand\kal{\mathcal}
\newcommand\fk{\mathfrak}
\newcommand\op{\oplus}
\newcommand\ot{\otimes}
\newcommand\ds{\displaystyle}
\newcommand\id{\indent}
\newcommand\nid{\noindent}
\newcommand\seq{\subseteq}
\newcommand\sneq{\subsetneq}
\newcommand\speq{\supseteq}
\newcommand\spneq{\supsetneq}
\newcommand\rar{\rightarrow}
\newcommand\Rar{\Rightarrow}
\newcommand\xrar{\xrightarrow}
\newcommand\lar{\leftarrow}
\newcommand\Lar{\Leftarrow}
\newcommand\xlar{\xleftarrow}
\newcommand\lrar{\leftrightarrow}
\newcommand\Lrar{\Leftrightarrow}
\newcommand\ideal{\trianglelefteq}
\newcommand\ai{\text{ a.\^{i}. }}
\newcommand\sk{(\textit{Schi\c{t}\u{a}}) }
\newcommand\rhup{\rightharpoonup}
\newcommand\rhdn{\rightharpoondown}
\newcommand\lhup{\leftharpoonup}
\newcommand\lhdn{\leftharpoondown}
\newcommand\vid{\emptyset}
\newcommand\wh{\widehat}
\newcommand\ol{\overline}
\newcommand\NN{\mathbb{N}}
\newcommand\RR{\mathbb{R}}
\newcommand\ZZ{\mathbb{Z}}
\newcommand\QQ{\mathbb{Q}}
\newcommand\CC{\mathbb{C}}

%% STILURILE MELE
\newtheoremstyle{dotless}{}{}{\itshape}{}{\bfseries}{:}{ }{}

\theoremstyle{dotless}
\newtheorem{theorem}{Teorem\u{a}}[section]
\newtheorem{proposition}{Propozi\c{t}ie}[section]
\newtheorem{lemma}{Lem\u{a}}[section]
\newtheorem{corollary}{Corolar}[section]
\newtheorem{exercise}{Exerci\c{t}iu}

\newtheoremstyle{dotlessdr}{}{}{}{}{\bfseries}{:}{ }{}
\theoremstyle{dotlessdr}
\newtheorem{example}{Exemplu}[section]
\newtheorem{definition}{Defini\c{t}ie}[section]
\newtheorem{remark}{Observa\c{t}ie}[section]




\begin{document}

\begin{center}
  {\large\textbf{
      Seminar 10 \\
      Extreme cu legături (condiționate)
    }}
\end{center}

\vspace{1cm}

\section{Extreme condiționate}
\label{sec:extreme-cond}

Atunci cînd domeniul de definiție al unei funcții de mai multe variabile conține,
la rîndul său anumite ecuații (numite, generic, \emph{legături}, problemele de
extrem se studiază folosind \textbf{metoda multiplicatorilor Lagrange}. Ea se
bazează pe următoarele concepte:

\begin{definition}\label{def:fc-scop}
  Fie $f(\pmb{x}, \pmb{y})$, cu $\pmb{x} \in \RR^n, \pmb{y} \in \RR^m$,
  iar $f : U \to \RR$ o funcție de $n + m$ variabile reale, cu valori reale, de
  clasă $\kal{C}^1$ pe un deschis $U \seq \RR^{n+m}$.

  Ea se numește \emph{funcție-scop (obiectiv)}. Presupunem că există $m$
  legături între variabilele $\pmb{x}, \pmb{y}$, adică $m$ relații de forma:
  \[
    g_1(\pmb{x}, \pmb{y}) = 0, \dots, g_m(\pmb{x}, \pmb{y}) = 0, \quad %
    g_i : U \to \RR,
  \]
  fiecare legătură fiind o funcție de clasă $\kal{C}^1(U)$.

  Fie $M = \{(x, y) \in U \mid g_i(\pmb{x}, \pmb{y}) = 0, 1 \leq i \leq m\}$
  mulțimea punctelor din $U$ care verifică legăturile.

  Se numește \emph{punct de extrem local al funcției $f$ cu legăturile $g_i$}
  orice punct $(\pmb{x}_0, \pmb{y}_0) \in M$, pentru care există o vecinătate
  $W \seq U$, astfel încît diferența $f(\pmb{x}, \pmb{y}) - f(\pmb{x}_0, \pmb{y}_0)$
  să aibă semn constant pentru orice $(\pmb{x}, \pmb{y}) \in M \cap W$.
\end{definition}

Teorema pe care se bazează metoda multiplicatorilor lui Lagrange este:

\begin{theorem}[J. L. Lagrange]\label{thm:lagr-mult}
  Cu notațiile și contextul de mai sus, presupunem că $(\pmb{x}_0, \pmb{y}_0)$
  este un punct de extrem local al lui $f$, cu legăturile de mai sus și că
  \[
    \det J_g(\pmb{x}_0, \pmb{y}_0) \stackrel{\text{not.}}{=\joinrel=\joinrel=} %
    \frac{D(g_1, \dots, g_m)}{D(y_1, \dots, y_m)} (\pmb{x}_0, \pmb{y}_0) \neq 0.
  \]

  Atunci există $m$ numere reale $\lambda_1, \dots, \lambda_m$, numite
  \emph{multiplicatori Lagrange} astfel încît, dacă definim funcția:
  \[
    F = f + \lambda_1 g_1 + \dots + \lambda_m g_m,
  \]
  punctul $(\pmb{x}_0, \pmb{y}_0)$ să verifice în mod necesar sistemul de
  $n + 2m$ ecuații:
  \[
    \frac{\partial F}{\partial x_j} = \frac{\partial F}{\partial y_k} = 0, %
    g_l = 0, \quad \forall 1 \leq j \leq n, 1 \leq k \leq m, 1 \leq l \leq m,
  \]
  cu $n + 2m$ necunoscute, $\pmb{\lambda}, \pmb{x}, \pmb{y}$.
\end{theorem}

\textbf{Exemplu:} Să se determine extremele locale ale funcției $f(x,y,z) = xyz$,
cu legătura $x + y + z = 1$.

\emph{Soluție:} Folosim metoda multiplicatorilor Lagrange. Definim funcțiile:
\begin{align*}
  g(x, y, z) &= x + y + z - 1 \\
  F(x, y, z) &= f + \lambda g = xyz + \lambda(x + y + z - 1).
\end{align*}

Extremele cerute verifică sistemul:
\[
  \begin{cases}
    \frac{\partial F}{\partial x} &= 0 \\
    \frac{\partial F}{\partial y} &= 0 \\
    \frac{\partial F}{\partial z} &= 0 \\
    g &= 0
  \end{cases} \Rightarrow
  \begin{cases}
    yz + \lambda &= 0 \\
    xz + \lambda &= 0 \\
    xy + \lambda &= 0 \\
    x + y + z - 1 &= 0
  \end{cases}
\]

Soluția sistemului este dată de:
\begin{align*}
  \lambda_1 &= -\frac{1}{9} \Rightarrow %
              (x_1, y_1, z_1) = \big(\frac{1}{3}, \frac{1}{3}, \frac{1}{3}\big) \\
  \lambda_2 &= 0 \Rightarrow
              \begin{cases}
                (x_2, y_2, z_2) &= (1, 0, 0) \\
                (x_3, y_3, z_3) &= (0, 1, 0) \\
                (x_4, y_4, z_4) &= (0, 0, 1)
              \end{cases}
\end{align*}

În continuare, studiem natura punctelor de extrem pentru funcția $F$, fie cu
matricea hessiană, fie cu diferențiala totală de ordin 2.

Găsim că $(x_1, y_1, z_1)$ este maxim local, iar celelalte puncte nu sînt
de extrem.

%%%%%%%%%%%%%%%%%%%%%%%%%%%%%%%%%%%%%%%%%%%%%%%%%%%%%%%%%%%%%%%%%%%%%%


\section{Exerciții}
\label{sec:exc}


1. Să se determine extremele funcțiilor $f$, cu legătura $g$ în cazurile:
\begin{enumerate}[(a)]
\item $f(x, y) = x^2 + y^2 - 6x - 2y + 1, g(x, y) = x + y - 2$;
\item $f(x, y) = x^2 + y^2 - 6x - 2y + 1, g(x, y) = x^2 + y^2 - 2x - 2y + 1$;
\item $f(x, y) = 3x + 4y, g(x, y) = x^2 + y^2 + 25$.
\end{enumerate}

\vspace{1cm}

2. Să se găsească punctul din planul $2x + y - z = 5$, situat la distanță minimă
față de origine.

\vspace{1cm}

3. Să se determine punctele cele mai depărtate de origine care se află pe
suprafața de ecuație $4x^2 + y^2 + z^2 - 8x - 4y + 4 = 0$.

\vspace{1cm}

4. Să se determine valorile extreme ale funcției:
\[
  f: \RR^3 \to \RR, \quad f(x, y, z) = 2x^2 + y^2 + 3z^2,
\]
pe mulțimea $\{(x, y, z) \in \RR^3 \mid x^2 + y^2 + z^2 = 1 \}$.

\vspace{1cm}

5. Să se determine valorile extreme ale produsului $xy$, cînd $x$ și $y$ sînt
coordonatele unui punct de pe elipsa de ecuație $x^2 + 2y^2 = 1$.
\end{document}
